
        \documentclass[spanish, 11pt]{exam}

        %These tell TeX which packages to use.
        \usepackage{array,epsfig}
        \usepackage{amsmath, textcomp}
        \usepackage{amsfonts}
        \usepackage{amssymb}
        \usepackage{amsxtra}
        \usepackage{amsthm}
        \usepackage{mathrsfs}
        \usepackage{color}
        \usepackage{multicol, xparse}
        \usepackage{verbatim}


        \usepackage[utf8]{inputenc}
        \usepackage[spanish]{babel}
        \usepackage{eurosym}

        \usepackage{graphicx}
        \graphicspath{{../img/}}
        \usepackage{pgf}
        
        \usepackage{pgfplots}
        \usetikzlibrary{math}
        \pgfplotsset{compat=1.15}
        \usepackage{xfp}

        \printanswers
        \nopointsinmargin
        \pointformat{}

        %Pagination stuff.
        %\setlength{\topmargin}{-.3 in}
        %\setlength{\oddsidemargin}{0in}
        %\setlength{\evensidemargin}{0in}
        %\setlength{\textheight}{9.in}
        %\setlength{\textwidth}{6.5in}
        %\pagestyle{empty}

        \let\multicolmulticols\multicols
        \let\endmulticolmulticols\endmulticols
        \RenewDocumentEnvironment{multicols}{mO{}}
         {%
          \ifnum#1=1
            #2%
          \else % More than 1 column
            \multicolmulticols{#1}[#2]
          \fi
         }
         {%
          \ifnum#1=1
          \else % More than 1 column
            \endmulticolmulticols
          \fi
         }
        \renewcommand{\solutiontitle}{\noindent\textbf{Sol:}\enspace}

        \newcommand{\samedir}{\mathbin{\!/\mkern-5mu/\!}}

        \newcommand{\class}{2º Bachillerato CCSS}
        \newcommand{\examdate}{\today}

        \newcommand{\tipo}{A}


        \newcommand{\timelimit}{50 minutos}



        \pagestyle{head}
        \firstpageheader{\includegraphics[width=0.2\columnwidth]{header_left}}{\textbf{Departamento de Matemáticas\linebreak \class}\linebreak \examnum}{\includegraphics[width=0.1\columnwidth]{header_right}}
        \runningheader{\class}{\examnum}{Página \thepage\ of \numpages}
        \runningheadrule

        \newcommand{\examnum}{Autoevaluación 2}
        \begin{document}
        \begin{questions}
        \question Se sabe que el gasto mensual en teléfono de los habitantes de una ciudad sigue una distribución normal
de desviación típica 6 euros. Se toma una muestra de 10 jóvenes y se obtienen los siguientes resultados (expresados 
en euros): \\ 24.5, 11, 16.5, 18.5, 21.5, 25, 6.5, 12,
10.5, 9.5  \\ Construya un intervalo de confianza de nivel 92 \% para la media del gasto mensual en 
teléfono de los habitantes 
de la ciudad. \begin{solution}   Calculamos el intervalo de confianza para la media, sabiendo que la media muestral es: 15.55, la desviación típica: 6, tamaño de la muestra: 10 y el grado de confianza: 92.0\%. \\ \\ Valor crítico: \\ $\alpha=1-0.92=0.08\to \frac{\alpha}{2}=0.04$ \\ \\ $P(Z>z_{\alpha/2})=0.04\to P(Z<z_{\alpha/2})=0.96 \to z_{\alpha/2} =1.7507$ \\ 
    \begin{tikzpicture}[scale=0.8]
    \pgfmathdeclarefunction{gauss}{2}{\pgfmathparse{1/(#2*sqrt(2*pi))*exp(-((x-#1)^2)/(2*#2^2))}}
    \tikzmath{\conf = 0.92; \crit= 1.7507; \a=0.04);}
    \begin{axis}[no markers, domain=-5:5, samples=100, axis lines=left, height=5cm, width=12cm, xtick={0,\crit}, ytick=\empty, xticklabels = {$0$, $z_{\frac{\alpha}{2}}=\crit$},enlargelimits=false, clip=false, axis on top]
      \addplot [fill=cyan!20, draw=none, domain=-\crit:\crit] {gauss(0,1)} \closedcycle;
      \addplot [very thick,cyan!50!black] {gauss(0,1)};
    \end{axis}
    \node[] at (5.2,1.5) {$\conf$};	
    \draw[-]   (\crit+6.5,1)node[right]{$\a$}  --  (\crit+5.6,0.1) ;
    \end{tikzpicture} \\
    Error cometido: \\ $E=z_{\alpha/2}\cdot \frac{\sigma}{\sqrt{n}} \to E=1.7507\cdot \frac{6}{3.1622776601683795}=3.3217$ \\ Por tanto el intervalo de confianza será: \\$\left(\overline{x} - E , \overline{x} + E \right)=\left(15.55 - 3.3217 , 15.55 + 3.3217 \right)=\left(12.2283, 18.8717 \right)$ \\  \\ 
    \begin{tikzpicture}[scale=0.4]
      \tikzmath{\a = -10; \b = 10; \aa = \a -1; \bb = \b + 1 ; \dist = \b - \a; \med = (\a + \b)/2;}
      \draw[very thick] (\a,0) -- (\b,0);
      \path [draw=black, fill=white] (\b,0) circle (2pt);
      \path [draw=black, fill=white] (\a,0.0) circle (2pt);
      \draw[latex-latex] (\a - 1.5,0) -- (\b + 1.5,0) ;
      \draw[shift={(\a,0)},color=black] (0pt,3pt) -- (0pt,-3pt);
      \draw[shift={(\a,0)},color=black] (0pt,0pt) -- (0pt,-3pt) node[below] {$12.2283$};
      \draw[shift={(\med,0)},color=black] (0pt,3pt) -- (0pt,-3pt);
      \draw[shift={(\med,0)},color=black] (0pt,0pt) -- (0pt,-3pt) node[below] {$15.55$};
      \draw[shift={(\b,0)},color=black] (0pt,3pt) -- (0pt,-3pt);
      \draw[shift={(\b,0)},color=black] (0pt,0pt) -- (0pt,-3pt) node[below] {$18.8717$};
      \draw[decorate,decoration={brace}, thick](\med,0.2)--(\b,0.2) node[above, midway] {$E=3.3217$}; 
    \end{tikzpicture} \\
       \end{solution}\question La duración de las bombillas de un fabricante es una variable aleatoria con distribución
normal de desviación típica 80 horas. \\ Queremos construir un intervalo de confianza al 97\% para la media de la duración de
las bombillas del fabricante, de forma que el error no sea mayor de 10 horas. ¿Qué tamaño de la muestra
debemos tomar? \begin{solution}   Calculamos el tamaño de la muestra para estimar la media con un error máximo: 10, si la desviación típica es: 80 y el grado de confianza: 97.0\%. \\ \\ Valor crítico: \\ $\alpha=1-0.97=0.03\to \frac{\alpha}{2}=0.015$ \\ \\ $P(Z>z_{\alpha/2})=0.015\to P(Z<z_{\alpha/2})=0.985 \to z_{\alpha/2} =2.1701$ \\ 
    \begin{tikzpicture}[scale=0.8]
    \pgfmathdeclarefunction{gauss}{2}{\pgfmathparse{1/(#2*sqrt(2*pi))*exp(-((x-#1)^2)/(2*#2^2))}}
    \tikzmath{\conf = 0.97; \crit= 2.1701; \a=0.015);}
    \begin{axis}[no markers, domain=-5:5, samples=100, axis lines=left, height=5cm, width=12cm, xtick={0,\crit}, ytick=\empty, xticklabels = {$0$, $z_{\frac{\alpha}{2}}=\crit$},enlargelimits=false, clip=false, axis on top]
      \addplot [fill=cyan!20, draw=none, domain=-\crit:\crit] {gauss(0,1)} \closedcycle;
      \addplot [very thick,cyan!50!black] {gauss(0,1)};
    \end{axis}
    \node[] at (5.2,1.5) {$\conf$};	
    \draw[-]   (\crit+6.5,1)node[right]{$\a$}  --  (\crit+5.6,0.1) ;
    \end{tikzpicture} \\
    Error cometido: \\ $E=z_{\alpha/2}\cdot \frac{\sigma}{\sqrt{n}} \to n \approx \frac{\sigma^2 \cdot z_{\alpha / 2}^2}{E^2}=301.39737664000006 \to n \geq302$ \\    \end{solution}\question La duración de las bombillas de un fabricante es una variable aleatoria con distribución
normal de desviación típica 80 horas. \\ Decidimos tomar un tamaño de la muestra igual a 150, comprobamos la duración de cada
bombilla y calculamos su promedio, que resulta ser 1053 horas. Calcular el intervalo de confianza al 97\%
para la media de la duración de las bombillas del fabricante. \begin{solution}   Calculamos el intervalo de confianza para la media, sabiendo que la media muestral es: 1053, la desviación típica: 80, tamaño de la muestra: 150 y el grado de confianza: 97.0\%. \\ \\ Valor crítico: \\ $\alpha=1-0.97=0.03\to \frac{\alpha}{2}=0.015$ \\ \\ $P(Z>z_{\alpha/2})=0.015\to P(Z<z_{\alpha/2})=0.985 \to z_{\alpha/2} =2.1701$ \\ 
    \begin{tikzpicture}[scale=0.8]
    \pgfmathdeclarefunction{gauss}{2}{\pgfmathparse{1/(#2*sqrt(2*pi))*exp(-((x-#1)^2)/(2*#2^2))}}
    \tikzmath{\conf = 0.97; \crit= 2.1701; \a=0.015);}
    \begin{axis}[no markers, domain=-5:5, samples=100, axis lines=left, height=5cm, width=12cm, xtick={0,\crit}, ytick=\empty, xticklabels = {$0$, $z_{\frac{\alpha}{2}}=\crit$},enlargelimits=false, clip=false, axis on top]
      \addplot [fill=cyan!20, draw=none, domain=-\crit:\crit] {gauss(0,1)} \closedcycle;
      \addplot [very thick,cyan!50!black] {gauss(0,1)};
    \end{axis}
    \node[] at (5.2,1.5) {$\conf$};	
    \draw[-]   (\crit+6.5,1)node[right]{$\a$}  --  (\crit+5.6,0.1) ;
    \end{tikzpicture} \\
    Error cometido: \\ $E=z_{\alpha/2}\cdot \frac{\sigma}{\sqrt{n}} \to E=2.1701\cdot \frac{80}{12.24744871391589}=14.175$ \\ Por tanto el intervalo de confianza será: \\$\left(\overline{x} - E , \overline{x} + E \right)=\left(1053 - 14.175 , 1053 + 14.175 \right)=\left(1038.825, 1067.175 \right)$ \\  \\ 
    \begin{tikzpicture}[scale=0.4]
      \tikzmath{\a = -10; \b = 10; \aa = \a -1; \bb = \b + 1 ; \dist = \b - \a; \med = (\a + \b)/2;}
      \draw[very thick] (\a,0) -- (\b,0);
      \path [draw=black, fill=white] (\b,0) circle (2pt);
      \path [draw=black, fill=white] (\a,0.0) circle (2pt);
      \draw[latex-latex] (\a - 1.5,0) -- (\b + 1.5,0) ;
      \draw[shift={(\a,0)},color=black] (0pt,3pt) -- (0pt,-3pt);
      \draw[shift={(\a,0)},color=black] (0pt,0pt) -- (0pt,-3pt) node[below] {$1038.825$};
      \draw[shift={(\med,0)},color=black] (0pt,3pt) -- (0pt,-3pt);
      \draw[shift={(\med,0)},color=black] (0pt,0pt) -- (0pt,-3pt) node[below] {$1053$};
      \draw[shift={(\b,0)},color=black] (0pt,3pt) -- (0pt,-3pt);
      \draw[shift={(\b,0)},color=black] (0pt,0pt) -- (0pt,-3pt) node[below] {$1067.175$};
      \draw[decorate,decoration={brace}, thick](\med,0.2)--(\b,0.2) node[above, midway] {$E=14.175$}; 
    \end{tikzpicture} \\
       \end{solution}\question Para optimizar los pedidos de las tallas de zapatos para un centro comercial, se
registraron los porcentajes del número de veces que fueron solicitadas las distintas tallas sobre una muestra de
1000 clientes. Si el porcentaje del número de veces para la talla media fueron del 25\%, calcule un intervalo de
confianza para tal porcentaje con un nivel de significación del 4\% \begin{solution}   Calculamos el intervalo de confianza para la proporción, si los aciertos son: 250.0, tamaño de la muestra: 1000 y el grado de confianza: 96.0\%. \\ \\ Valor crítico: \\ $\alpha=1-0.96=0.04\to \frac{\alpha}{2}=0.02$ \\ \\ $P(Z>z_{\alpha/2})=0.02\to P(Z<z_{\alpha/2})=0.98 \to z_{\alpha/2} =2.0537$ \\ 
    \begin{tikzpicture}[scale=0.8]
    \pgfmathdeclarefunction{gauss}{2}{\pgfmathparse{1/(#2*sqrt(2*pi))*exp(-((x-#1)^2)/(2*#2^2))}}
    \tikzmath{\conf = 0.96; \crit= 2.0537; \a=0.02);}
    \begin{axis}[no markers, domain=-5:5, samples=100, axis lines=left, height=5cm, width=12cm, xtick={0,\crit}, ytick=\empty, xticklabels = {$0$, $z_{\frac{\alpha}{2}}=\crit$},enlargelimits=false, clip=false, axis on top]
      \addplot [fill=cyan!20, draw=none, domain=-\crit:\crit] {gauss(0,1)} \closedcycle;
      \addplot [very thick,cyan!50!black] {gauss(0,1)};
    \end{axis}
    \node[] at (5.2,1.5) {$\conf$};	
    \draw[-]   (\crit+6.5,1)node[right]{$\a$}  --  (\crit+5.6,0.1) ;
    \end{tikzpicture} \\
    Error cometido: \\ $E=z_{\alpha / 2}\cdot \sqrt{\frac{\widehat{p}\cdot\left(1-\widehat{p} \right)}{n}} \to E=2.0537\cdot \sqrt{\frac{0.25\cdot0.75}{1000}}=0.0281$ \\ Por tanto el intervalo de confianza será: \\$\left(\widehat{p} - E , \widehat{p} + E \right)=\left(0.25 - 0.0281 , 0.25 + 0.0281 \right)=\left(0.2219, 0.2781 \right)$ \\  \\ 
    \begin{tikzpicture}[scale=0.4]
      \tikzmath{\a = -10; \b = 10; \aa = \a -1; \bb = \b + 1 ; \dist = \b - \a; \med = (\a + \b)/2;}
      \draw[very thick] (\a,0) -- (\b,0);
      \path [draw=black, fill=white] (\b,0) circle (2pt);
      \path [draw=black, fill=white] (\a,0.0) circle (2pt);
      \draw[latex-latex] (\a - 1.5,0) -- (\b + 1.5,0) ;
      \draw[shift={(\a,0)},color=black] (0pt,3pt) -- (0pt,-3pt);
      \draw[shift={(\a,0)},color=black] (0pt,0pt) -- (0pt,-3pt) node[below] {$0.2219$};
      \draw[shift={(\med,0)},color=black] (0pt,3pt) -- (0pt,-3pt);
      \draw[shift={(\med,0)},color=black] (0pt,0pt) -- (0pt,-3pt) node[below] {$0.25$};
      \draw[shift={(\b,0)},color=black] (0pt,3pt) -- (0pt,-3pt);
      \draw[shift={(\b,0)},color=black] (0pt,0pt) -- (0pt,-3pt) node[below] {$0.2781$};
      \draw[decorate,decoration={brace}, thick](\med,0.2)--(\b,0.2) node[above, midway] {$E=0.0281$}; 
    \end{tikzpicture} \\
       \end{solution}\question Para estimar la proporción de las familias de una determinada ciudad que tienen smart tv, 
se quiere
utilizar una muestra aleatoria de medida n. Calcula el valor mínimo de n para garantizar que, a un nivel de
confianza del 92\%, el error en la estimación sea menor que 0,05. (PISTA: Como se desconoce la proporción, 
se ha de tomar el caso más desfavorable que será 0.5) \begin{solution}   Calculamos el tamaño de la muestra para estimar la proporcion con un error máximo: 0.05, si la proporción muestral es: 0.5 y el grado de confianza: 92.0\%. \\ \\ Valor crítico: \\ $\alpha=1-0.92=0.08\to \frac{\alpha}{2}=0.04$ \\ \\ $P(Z>z_{\alpha/2})=0.04\to P(Z<z_{\alpha/2})=0.96 \to z_{\alpha/2} =1.7507$ \\ 
    \begin{tikzpicture}[scale=0.8]
    \pgfmathdeclarefunction{gauss}{2}{\pgfmathparse{1/(#2*sqrt(2*pi))*exp(-((x-#1)^2)/(2*#2^2))}}
    \tikzmath{\conf = 0.92; \crit= 1.7507; \a=0.04);}
    \begin{axis}[no markers, domain=-5:5, samples=100, axis lines=left, height=5cm, width=12cm, xtick={0,\crit}, ytick=\empty, xticklabels = {$0$, $z_{\frac{\alpha}{2}}=\crit$},enlargelimits=false, clip=false, axis on top]
      \addplot [fill=cyan!20, draw=none, domain=-\crit:\crit] {gauss(0,1)} \closedcycle;
      \addplot [very thick,cyan!50!black] {gauss(0,1)};
    \end{axis}
    \node[] at (5.2,1.5) {$\conf$};	
    \draw[-]   (\crit+6.5,1)node[right]{$\a$}  --  (\crit+5.6,0.1) ;
    \end{tikzpicture} \\
    Error cometido: \\ $E=z_{\alpha / 2}\cdot \sqrt{\frac{\widehat{p}\cdot\left(1-\widehat{p} \right)}{n}} \to n =\frac{\widehat{p}\cdot\left(1-\widehat{p} \right)\cdot z_{\alpha / 2}^2}{E^2} \to n \geq307$ \\    \end{solution}\question La cantidad de horas semanales que los habitantes de Utebo emplean en ver series de TV
se puede aproximar
por una distribución normal con una desviación típica de 0,64. Se toma una muestra aleatoria simple y se
obtienen los siguientes datos: \\
6.9, 7.3, 7.6, 6.6, 6.5, 7.1, 6.2, 6.9, 7.8, 6.7, 7.0, 6.5, 5.5, 7.2, 7.6, 5.8 \\

Determinar el nivel de confianza para el cual el intervalo de confianza para la media de
horas semanales que se ven series de TV es (6.65 , 7). Detallar los pasos realizados para obtener los resultados. \begin{solution}   La media muestral es: 6.825 \\ Del intervalo obtenemos la media muestral: 6.825 y el error máximo: $6.825-6.65=0.175$ \\ Calculamos el grado de confianza para un error máximo: 0.175, siendo el tamaño muestral: 16 y la desviación típica: 0.64. \\ \\  \\ Como $E=z_{\alpha / 2}\cdot \frac{\sigma}{\sqrt{n}} \Rightarrow z_{\alpha / 2} =\frac{E \cdot \sqrt{n}}{\sigma}\to z_{\alpha / 2}=1.09$ \\ $P(Z>z_{\alpha / 2})=P(Z>1.09)=\frac{\alpha}{2}=0.1379$ \\ $\alpha=0.2757 \to \quad confianza=1 - \alpha=0.7243 \to 72.43 \%$    \end{solution}
    \end{questions}
    \end{document}
    